\subsection{Entorno Experimental, Infraestructura y Gestión de Cuotas}
\label{subsec:entorno_experimental_infraestructura}

Esta subsección describe la plataforma de cómputo, las librerías del entorno de ejecución y las políticas de control de cuotas y memoria del prototipo.

\subsubsection{Plataforma de Ejecución y Pila Tecnológica}
Los experimentos y el desarrollo del prototipo se ejecutaron sobre una estación de trabajo con procesador Apple Silicon M3 (arquitectura ARM64), $16\text{ GB}$ de memoria unificada y sistema operativo macOS, utilizando el entorno de ejecución Python 3.13 (\texttt{3.13.12}). La pila de software científico y librerías especializadas comprende:
\begin{itemize}
    \item \textbf{Procesamiento Audiovisual:} \texttt{OpenCV} (\texttt{opencv-python}) para la decodificación de metraje, cálculo matricial del índice SSIM, muestreo por entropía visual y compresión JPEG; \texttt{ffmpeg} para el remuestreo de audio; y \texttt{openai-whisper} (modelo \texttt{base}) para el reconocimiento automático del habla (ASR) ejecutado localmente en CPU.
    \item \textbf{Modelos Fundacionales y Motores de Vectorización (Cloud y Local):} SDK oficial \texttt{google-genai} para el acceso remoto a \textit{Gemini Embedding 2} (\texttt{gemini-embedding-2-preview}) y el árbitro bimodal \textit{Gemini 3.5 Flash Lite} (\texttt{gemini-3.5-flash-lite}); junto con \texttt{PyTorch} (\texttt{2.14.1}, acelerador Metal Performance Shaders \texttt{mps}), \texttt{transformers} (\texttt{5.22.2}), \texttt{sentence-transformers[image]} (\texttt{6.3.0}) y \texttt{torchvision} (\texttt{0.29.2}) para inferencia local de \textit{Qwen3-VL-Embedding-2B} y \textit{EmbeddingGemma 2} en precisión \texttt{bfloat16} sobre \texttt{mps} (y \texttt{float32} en CPU).
    \item \textbf{Indexación Vectorial:} \texttt{qdrant-client} en modo local persistente en disco (\texttt{./data/vector\_db}), evitando la sobrecarga de latencia de red y garantizando aislamiento experimental entre configuraciones mediante colecciones independientes.
    \item \textbf{Cálculo Numérico, Estadístico y Evaluación:} \texttt{NumPy} para el truncamiento dimensional y normalización $L_2$ bajo el estándar MRL, \texttt{SciPy} (\texttt{1.18.1}) para las pruebas de significación pareada (\textit{McNemar} y \textit{Wilcoxon}), y \texttt{pytest} como motor de ejecución para las suites de prueba unitarias y de extremo a extremo.
\end{itemize}

\subsubsection{Gestión de Cuotas, Concurrencia y Presupuesto de Memoria Unificada}
La capa gratuita de la API de Google Gemini impone restricciones severas que condicionan el diseño de cualquier sistema experimental en tiempo real: un máximo de $100$ peticiones por minuto (RPM), $30.000$ tokens por minuto (TPM) y $1.000$ peticiones diarias (RPD) \cite{gemini_emb2_docs}. En paralelo, la ejecución local de modelos multimodales de hasta $2\text{B}$ parámetros dentro de $16\text{ GB}$ de RAM unificada requiere evitar la fragmentación de caché en GPU (\texttt{mps}) y las colisiones concurrentes de descriptores de comandos.

Para operar de manera robusta y continua sin interrupciones por denegación de servicio (código HTTP 429 / \texttt{RESOURCE\_EXHAUSTED}) ni desbordamiento de memoria, se implementaron cuatro salvaguardas arquitectónicas:
\begin{enumerate}
    \item \textbf{Reintentos Exponenciales con Dispersión (\textit{Exponential Backoff with Jitter}):} Toda llamada remota a la API se canaliza a través de la función utilitaria \texttt{call\_with\_retry}. Ante errores transitorios de cuota o sobrecarga de red, el cliente aplica una pausa adaptativa con retardo base de $4,0\text{ s}$ para vectorización y $1,5\text{ s}$ para reordenación, incrementándose exponencialmente con un factor de aleatoriedad hasta un máximo de 4 reintentos antes de elevar una excepción controlada (\texttt{VectorizationError}).
    \item \textbf{Cortafuegos de Cuota por Vídeo (\texttt{QuotaSafetyError}):} Antes de vectorizar, el orquestador verifica que el total de fragmentos no supere $\text{MAX\_CHUNKS\_PER\_VIDEO} = 800$. Si un vídeo supera dicho umbral, aborta preventivamente antes de emitir llamadas a la API, evitando agotar la cuota diaria.
    \item \textbf{Semáforo Adaptativo y Cerrojo por Modelo:} En \texttt{process\_\allowbreak video\_\allowbreak pipeline}, las llamadas remotas a Gemini emplean un máximo de $3$ peticiones concurrentes (\texttt{max\_\allowbreak concurrent\_\allowbreak requests = 3}), mientras que los motores locales fuerzan concurrencia estricta de $1$ (\texttt{Semaphore(1)}) y serializan las pasadas \texttt{model.encode} mediante un \texttt{threading.Lock} compartido por \texttt{model\_id} en \texttt{local\_\allowbreak vectorizer.py}.
    \item \textbf{Descarga Secuencial de Modelos (\textit{Single-Resident Lifecycle}):} Antes de cargar un vectorizador local tras la transcripción ASR, el sistema invoca \texttt{unload\_whisper\_models()}; de manera simétrica, la factoría \texttt{factory.py} mantiene un único modelo local residente y libera el anterior mediante \texttt{unload\_active\_vectorizer()} (\texttt{gc.collect()} seguido de \texttt{torch.mps.empty\_cache()}).
\end{enumerate}

\subsubsection{Configuración y Aislamiento de Colecciones en Qdrant}
La base vectorial local almacena los fragmentos indexados en colecciones aisladas generadas por \texttt{collection\_name\_for(backend, dim)}. Para preservar el índice histórico de referencia sin consumir cuota diaria adicional, la configuración predeterminada (\texttt{gemini} en $d=3072$) conserva la colección \texttt{video\_chunks}, mientras que cualquier otra combinación de motor y dimensionalidad se aísla bajo el esquema determinista \texttt{video\_chunks\_\_<backend>\_\_<dim>} (cubriendo \texttt{gemini} en $1536/768$, \texttt{qwen3-vl-2b} en $2048/1536/768$ y \texttt{gemma2} en $768/512/256$). Cada colección se configura con métrica de distancia coseno (\texttt{Distance.COSINE}) y almacena un \textit{payload} documental uniforme (nombre de vídeo, \texttt{chunk\_index}, intervalo $[t_{\text{start}}, t_{\text{end}}]$ y transcripción), habilitando la conmutación instantánea de motores en la interfaz sin reindexar.
